% =====================================================================
%  รายงานโครงงาน: การตรวจจับ Wash Trading บน Poly\-market
%  เวอร์ชัน LaTeX ให้ตรงกับไฟล์ Word (ฟอนต์ TH SarabunPSK, ตัวหนา, ดำล้วน)
%
%  คอมไพล์ด้วย XeLaTeX สองรอบ (เพื่อให้เลขหน้าในสารบัญถูกต้อง):
%      xelatex WashTrading_Report_TH.tex
%      xelatex WashTrading_Report_TH.tex
%  ต้องติดตั้งฟอนต์ TH SarabunPSK ในเครื่อง (เครื่องที่เปิดไฟล์ Word ก็มีอยู่แล้ว)
% =====================================================================
\documentclass[12pt]{article}

\usepackage[letterpaper,left=1.5in,right=1in,top=1.5in,bottom=1in]{geometry}
\usepackage{fontspec}
\usepackage{polyglossia}
\setmainlanguage{thai}
\setotherlanguage{english}

% ฟอนต์หลักทั้งเอกสาร — เปลี่ยนชื่อได้ที่บรรทัดนี้บรรทัดเดียว
\newcommand{\REPORTFONT}{TH Sarabun PSK}
\setmainfont{\REPORTFONT}
% บอก polyglossia ให้ใช้ฟอนต์นี้กับภาษาไทยโดยตรง (กันคำเตือน/ข้อผิดพลาดเรื่อง Thai script)
\newfontfamily\thaifont{\REPORTFONT}

\usepackage[table]{xcolor}
\usepackage{array}
\usepackage{enumitem}
\usepackage{setspace}
\usepackage{fancyhdr}
\usepackage{graphicx}
\usepackage{url}
\urlstyle{same}   % ลิงก์ใช้ฟอนต์เดียวกับเนื้อหา (ไม่ใช่ typewriter)

\definecolor{hdrgray}{HTML}{D9D9D9}

% ฟอนต์สำรองเฉพาะอักขระที่ TH Sarabun PSK ไม่มี (เช่น θ) ปรับขนาดให้เท่าตัวอักษรหลัก
\newfontfamily\symfont{DejaVu Sans}[Scale=MatchLowercase]
\newcommand{\gtheta}{{\symfont\bfseries θ}}

% ขนาดฐาน 16 pt และตัวหนา–ดำล้วนทั้งเอกสาร
\newcommand{\basesize}{\fontsize{16}{24}\selectfont}
\AtBeginDocument{\basesize\bfseries\color{black}}
\setstretch{1.3}
\setlength{\parindent}{0.5in}
\setlength{\parskip}{0.35em}

% ยังจัดชิดขอบสองข้าง (justified) แต่คุมช่องว่างไม่ให้ยืดมากเกินไป
% แทนการใช้ \sloppy (ซึ่งยอมให้บรรทัดหลวมมาก) ด้วยค่าที่พอดีกว่า
\tolerance=2000
\emergencystretch=2em

% เลขหน้าอยู่กึ่งกลางด้านล่าง
\pagestyle{fancy}
\fancyhf{}
\renewcommand{\headrulewidth}{0pt}
\cfoot{\basesize\bfseries\thepage}

% ---------- หัวข้อระดับบท (จัดกึ่งกลาง ขึ้นหน้าใหม่) ----------
\newcommand{\chap}[3]{% #1 = เลขบท, #2 = ชื่อบท, #3 = label
  \clearpage\label{#3}%
  \begin{center}{\fontsize{20}{28}\selectfont\bfseries บทที่ #1\par\vspace{2pt}#2\par}\end{center}%
  \vspace{0.8em}\par}
% หัวข้อไม่มีเลขบท (บทคัดย่อ/คำนำ/สารบัญ/บรรณานุกรม)
\newcommand{\plainchap}[2]{% #1 = ชื่อ, #2 = label
  \clearpage\label{#2}%
  \begin{center}{\fontsize{20}{28}\selectfont\bfseries #1\par}\end{center}%
  \vspace{0.8em}\par}
% หัวข้อย่อย (ชิดซ้าย) — ช่องว่างบน/ล่างคงที่ทุกหัวข้อ ไม่ขึ้นกับว่าก่อน/หลังเป็นย่อหน้าหรือรายการ
% \addvspace คุมช่องว่างด้านบนให้ "เท่ากันเสมอ" (ดูดกลืนช่องว่างจากย่อหน้า/รายการ/หัวบทที่มาก่อน)
\newcommand{\sect}[2]{% #1 = ข้อความ, #2 = label
  \par\addvspace{16pt}\label{#2}%
  \noindent{\fontsize{17}{23}\selectfont\bfseries #1}%
  \par\nobreak\vspace{4pt}}
% หัวข้อย่อยลึกอีกชั้น (เช่น 3.3.1) — ไม่ลงสารบัญ
\newcommand{\subsect}[1]{\par\vspace{0.3em}\noindent\hspace{0.3in}{\bfseries #1}\par\vspace{0.1em}}

% ---------- บรรทัดในสารบัญ: ข้อความ ... เลขหน้า ----------
\newcommand{\tocI}[2]{\par\vspace{4pt}{\fontsize{16}{20}\selectfont\noindent\mbox{#1}\nobreak\dotfill\pageref{#2}\par}}
\newcommand{\tocII}[2]{\par\vspace{2pt}{\fontsize{16}{20}\selectfont\noindent\hspace*{0.35in}\mbox{#1}\nobreak\dotfill\pageref{#2}\par}}

% ---------- รายการ ----------
% topsep/partopsep ตั้งให้ช่องว่างเหนือรายการเท่ากับ \parskip ของย่อหน้า
% เพื่อให้ช่องว่างใต้หัวข้อ "เท่ากัน" ไม่ว่าจะตามด้วยย่อหน้าหรือรายการ
\newlist{numlist}{enumerate}{1}
\setlist[numlist]{label=\arabic*), leftmargin=*, labelsep=0.5em, topsep=0.35em, partopsep=0pt, parsep=0pt, itemsep=3pt}
\newlist{dotlist}{itemize}{1}
\setlist[dotlist]{label=\textbullet, leftmargin=*, labelsep=0.5em, topsep=0.35em, partopsep=0pt, parsep=0pt, itemsep=3pt}

% ---------- ตาราง 2 คอลัมน์ / 3 คอลัมน์ ----------
\newcolumntype{L}[1]{>{\raggedright\arraybackslash}p{#1}}

\begin{document}

% =====================================================================
% ปก
% =====================================================================
\thispagestyle{empty}
\begin{center}
  \vspace*{1.2cm}
  {\fontsize{26}{32}\selectfont\bfseries รายงานกลุ่ม}\par\vspace{0.6cm}
  {\fontsize{16}{22}\selectfont วิชา Analytics on Complex Data (01204465)}\par\vspace{0.2cm}
  {\fontsize{16}{22}\selectfont Assignment \#2 — Group Assignment}\par\vspace{1.4cm}

  {\fontsize{20}{28}\selectfont\bfseries การตรวจจับการซื้อขายสมรู้ร่วมคิด (Wash Trading)\\
   บนแพลตฟอร์ม Poly\-market ด้วยการวิเคราะห์ข้อมูลเชิงเครือข่าย}\par\vspace{0.5cm}
  {\fontsize{16}{22}\selectfont\textit{Detecting Wash Trading on Poly\-market using Network-based Analytics}}\par\vspace{1.6cm}

  {\fontsize{16}{22}\selectfont\bfseries จัดทำโดย กลุ่มที่ 3}\par\vspace{2.0cm}

  {\fontsize{16}{22}\selectfont\bfseries คณะวิศวกรรมศาสตร์ มหาวิทยาลัยเกษตรศาสตร์}\par\vspace{0.2cm}
  {\fontsize{16}{22}\selectfont ภาคต้น ปีการศึกษา 2567}\par
\end{center}

% =====================================================================
% บทคัดย่อ
% =====================================================================
\plainchap{บทคัดย่อ}{abstract}
รายงานนี้ศึกษาการตรวจจับการซื้อขายสมรู้ร่วมคิด (wash trading) บนแพลตฟอร์มตลาดทำนายผล Poly\-market ซึ่งเป็นพฤติกรรมที่ผู้ใช้ซื้อขายวนกันเองเพื่อปั่นปริมาณการซื้อขาย (volume) ให้ดูสูงเกินจริง ข้อมูลที่ใช้เป็นข้อมูลธุรกรรมบนบล็อกเชน (on-chain) ที่มีลักษณะเป็นข้อมูลเชิงเครือข่าย (graph/network data) โดยมองกระเป๋าเงิน (wallet) เป็นโหนด และการซื้อขายระหว่างกระเป๋าเป็นเส้นเชื่อม

คณะผู้จัดทำได้นำ Algorithm 1 จากงานวิจัยของ Sirolly และคณะ (2025) มาพัฒนาใหม่เป็นวิธีให้คะแนนความน่าสงสัยแก่ทุกกระเป๋า แล้วทดสอบสมมติฐานเบื้องต้น (premise test) โดยเปรียบเทียบช่วงที่งานวิจัยระบุว่ามี wash สูง (ธันวาคม 2567) กับช่วงที่ต่ำ (มิถุนายน 2568) ผลการวิเคราะห์ครบ 18 สัปดาห์ (ประมาณ 29.38 ล้านรายการ) พบสัดส่วนปริมาณการซื้อขายที่ถูกติดธงเท่ากับ 49.00\% เทียบกับ 2.50\% ที่ค่าเกณฑ์ \gtheta\ = 0.90 ความแตกต่างนี้คงอยู่ที่ทุกค่าเกณฑ์ที่ทดลอง

ทั้งนี้ ผลที่ได้เป็นการตรวจจับ ``รูปแบบที่น่าสงสัย'' ไม่ใช่การยืนยันว่าทุกรายการที่ถูกติดธงเป็น wash trading จริง เนื่องจากไม่มีข้อมูลเฉลย (ground truth) ที่แท้จริง คณะผู้จัดทำจึงได้ประเมินเพิ่มเติมด้วยการฉีดรูปแบบที่รู้คำตอบ (synthetic injection) และรายงานข้อจำกัดไว้อย่างตรงไปตรงมา

\vspace{0.4em}\noindent\textbf{คำสำคัญ:} wash trading, Poly\-market, การวิเคราะห์เครือข่าย, ข้อมูลเชิงซ้อน, การตรวจจับความผิดปกติ

% =====================================================================
% คำนำ
% =====================================================================
\plainchap{คำนำ}{preface}
รายงานฉบับนี้เป็นส่วนหนึ่งของวิชา Analytics on Complex Data (01204465) โดยมีวัตถุประสงค์เพื่อแสดงให้เห็นว่ากลุ่มของเราสามารถสำรวจ เตรียม วิเคราะห์ และประเมินผลข้อมูลเชิงซ้อน (complex data) ได้ ผ่านกรณีศึกษาการตรวจจับ wash trading บนแพลตฟอร์ม Poly\-market ซึ่งใช้ข้อมูลเชิงเครือข่ายและเทคนิคการให้คะแนนแบบเครือข่ายที่ไม่ได้สอนในชั้นเรียน

เนื้อหาในรายงานครอบคลุมตั้งแต่ที่มาและความสำคัญของปัญหา ทฤษฎีและงานวิจัยที่เกี่ยวข้อง วิธีดำเนินการตั้งแต่การเตรียมข้อมูลจนถึงการพัฒนาอัลกอริทึม ผลการศึกษา ตลอดจนข้อจำกัดและข้อเสนอแนะ คณะผู้จัดทำหวังว่ารายงานนี้จะเป็นประโยชน์ต่อผู้ที่สนใจ และขอขอบคุณอาจารย์ผู้สอนที่ให้คำแนะนำมา ณ โอกาสนี้

\vspace{0.6em}\hfill\textbf{คณะผู้จัดทำ}

% =====================================================================
% สารบัญ (เลขหน้าอ้างอิงอัตโนมัติด้วย \pageref)
% =====================================================================
\clearpage
\begin{center}{\fontsize{20}{28}\selectfont\bfseries สารบัญ}\end{center}
\vspace{0.6em}
{\basesize\setstretch{1.25}
\tocI{บทคัดย่อ}{abstract}
\tocI{คำนำ}{preface}
\tocI{บทที่ 1\quad บทนำ}{ch1}
\tocII{1.1 ความเป็นมาและความสำคัญของปัญหา}{s11}
\tocII{1.2 วัตถุประสงค์}{s12}
\tocII{1.3 ขอบเขตของการศึกษา}{s13}
\tocII{1.4 ประโยชน์ที่คาดว่าจะได้รับ}{s14}
\tocI{บทที่ 2\quad ทฤษฎีและงานวิจัยที่เกี่ยวข้อง}{ch2}
\tocII{2.1 Wash Trading และตลาดทำนายผล}{s21}
\tocII{2.2 ข้อมูลเชิงเครือข่าย (Graph / Network Data)}{s22}
\tocII{2.3 Algorithm 1 (Sirolly et al., 2025)}{s23}
\tocII{2.4 เครื่องมือและไลบรารีที่เกี่ยวข้อง}{s24}
\tocI{บทที่ 3\quad วิธีดำเนินการศึกษา}{ch3}
\tocII{3.1 แหล่งข้อมูลและการเก็บรวบรวม}{s31}
\tocII{3.2 การเตรียมข้อมูล (Pre-processing)}{s32}
\tocII{3.3 การพัฒนา Algorithm 1}{s33}
\tocII{3.4 การทดสอบสมมติฐานและการประเมินผล}{s34}
\tocII{3.5 การวิเคราะห์ครบ 18 สัปดาห์และการสำรวจข้อมูล (EDA)}{s35}
\tocI{บทที่ 4\quad ผลการศึกษา}{ch4}
\tocII{4.1 ผลการทดสอบสมมติฐาน (Baseline)}{s41}
\tocII{4.2 ผลการวิเคราะห์ครบ 18 สัปดาห์}{s42}
\tocII{4.3 ความไวต่อค่าเกณฑ์ (\gtheta-sensitivity)}{s43}
\tocII{4.4 การสำรวจข้อมูลกระเป๋าที่ถูกติดธง}{s44}
\tocII{4.5 การประเมินด้วยข้อมูลที่รู้คำตอบ (Synthetic Evaluation)}{s45}
\tocII{4.6 กรณีที่ต้องตีความอย่างระมัดระวัง}{s46}
\tocI{บทที่ 5\quad สรุปผลและข้อเสนอแนะ}{ch5}
\tocII{5.1 สรุปผลการศึกษา}{s51}
\tocII{5.2 ข้อจำกัด}{s52}
\tocII{5.3 ข้อเสนอแนะ}{s53}
\tocI{บรรณานุกรม}{refs}
}

% =====================================================================
% บทที่ 1
% =====================================================================
\chap{1}{บทนำ}{ch1}
\sect{1.1 ความเป็นมาและความสำคัญของปัญหา}{s11}
Poly\-market เป็นแพลตฟอร์มตลาดทำนายผล (prediction market) แบบกระจายศูนย์บนบล็อกเชน ผู้ใช้สามารถซื้อขาย ``หุ้นผลลัพธ์'' ของเหตุการณ์ต่าง ๆ เช่น ผลการเลือกตั้งหรือผลกีฬา ปริมาณการซื้อขาย (trading volume) เป็นตัวชี้วัดสำคัญที่สะท้อนความนิยมและสภาพคล่องของแต่ละตลาด

อย่างไรก็ตาม ปริมาณการซื้อขายอาจถูกปั่นให้สูงเกินจริงด้วยพฤติกรรมที่เรียกว่า wash trading ซึ่งเป็นการที่กลุ่มกระเป๋าเงินซื้อขายวนไปมากันเองโดยไม่ได้เปลี่ยนความเสี่ยงที่แท้จริง เพื่อสร้างภาพลวงว่าตลาดมีความคึกคัก พฤติกรรมนี้ทำให้ข้อมูลปริมาณการซื้อขายบิดเบือนและอาจหลอกลวงผู้ใช้รายอื่น การตรวจจับ wash trading จึงเป็นปัญหาสำคัญทั้งในด้านความโปร่งใสของตลาดและการวิเคราะห์ข้อมูล

ข้อมูลธุรกรรมบนบล็อกเชนของ Poly\-market เป็นข้อมูลเชิงซ้อนประเภทข้อมูลเชิงเครือข่าย (graph/network data) ที่ไม่สามารถวิเคราะห์ด้วยตารางแบนธรรมดาได้ เนื่องจากความสัมพันธ์ระหว่างกระเป๋า (ใครเทรดกับใคร และกระจุกตัวกับคู่เทรดมากน้อยเพียงใด) คือหัวใจของการแยกแยะพฤติกรรมปกติออกจาก wash trading รายงานนี้จึงเลือกใช้เทคนิคการวิเคราะห์เชิงเครือข่ายเพื่อแก้ปัญหาดังกล่าว

\sect{1.2 วัตถุประสงค์}{s12}
\begin{numlist}
\item เพื่อพัฒนาวิธีให้คะแนนความน่าสงสัยในการทำ wash trading แก่กระเป๋าเงินบน Poly\-market ด้วยการวิเคราะห์เชิงเครือข่าย
\item เพื่อทดสอบสมมติฐานว่าวิธีที่พัฒนาสามารถแยกช่วงเวลาที่มี wash สูงออกจากช่วงที่มี wash ต่ำได้อย่างชัดเจน
\item เพื่อประเมินคุณภาพของตัวตรวจจับด้วยข้อมูลที่รู้คำตอบ (synthetic evaluation) และวิเคราะห์ความไวต่อพารามิเตอร์
\item เพื่อฝึกการจัดการข้อมูลเชิงซ้อนครบวงจร ตั้งแต่การเตรียมข้อมูล การสำรวจ การวิเคราะห์ จนถึงการประเมินผล
\end{numlist}

\sect{1.3 ขอบเขตของการศึกษา}{s13}
\begin{dotlist}
\item ข้อมูล: ชุดข้อมูล on-chain ของ Poly\-market (moose-code/polymarket-onchain-v1) บน Hugging Face สัญญาอนุญาต CC-BY-4.0
\item ช่วงเวลา: พฤศจิกายน–ธันวาคม 2567 และ พฤษภาคม–มิถุนายน 2568 รวม 18 สัปดาห์ (122 วันเป้าหมาย ประมาณ 29.38 ล้านรายการ)
\item เทคนิค: การให้คะแนนเชิงเครือข่ายตาม Algorithm 1 (Sirolly et al., 2025) ร่วมกับการสำรวจข้อมูล (EDA) และการประเมินแบบฉีดรูปแบบที่รู้คำตอบ
\item ขอบเขตที่ไม่รวม: การยืนยันเจตนาสมรู้ร่วมคิดทางกฎหมาย และการพิสูจน์ ground truth ของแต่ละรายการ
\end{dotlist}

\sect{1.4 ประโยชน์ที่คาดว่าจะได้รับ}{s14}
\begin{numlist}
\item ได้เครื่องมือต้นแบบสำหรับตรวจจับรูปแบบ wash trading บนตลาดทำนายผลแบบกระจายศูนย์
\item เข้าใจแนวทางการจัดการและวิเคราะห์ข้อมูลเชิงเครือข่ายขนาดใหญ่ในทางปฏิบัติ
\item เข้าใจข้อจำกัดของการตรวจจับความผิดปกติเมื่อไม่มีข้อมูลเฉลย และแนวทางการประเมินที่โปร่งใส
\end{numlist}

% =====================================================================
% บทที่ 2
% =====================================================================
\chap{2}{ทฤษฎีและงานวิจัยที่เกี่ยวข้อง}{ch2}
\sect{2.1 Wash Trading และตลาดทำนายผล}{s21}
Wash trading คือการซื้อขายสินทรัพย์โดยที่ผู้ซื้อและผู้ขายเป็นกลุ่มเดียวกันหรือสมรู้ร่วมคิดกัน ทำให้ความเป็นเจ้าของที่แท้จริงไม่เปลี่ยนแปลง แต่สร้างปริมาณการซื้อขายปลอมขึ้นมา บนตลาดทำนายผลอย่าง Poly\-market พฤติกรรมนี้มักปรากฏเป็นการเปิดและปิดสถานะ (position) ซ้ำ ๆ ระหว่างกระเป๋าไม่กี่ใบ

\sect{2.2 ข้อมูลเชิงเครือข่าย (Graph / Network Data)}{s22}
ข้อมูลเชิงเครือข่ายเป็นข้อมูลเชิงซ้อนประเภทหนึ่งที่โครงสร้างอยู่ในรูปของโหนดและเส้นเชื่อม ไม่สามารถแทนด้วยตารางแบนได้โดยตรง ในงานนี้กระเป๋าเงินเป็นโหนด และการซื้อขายระหว่างกระเป๋าเป็นเส้นเชื่อมที่มีน้ำหนักตามปริมาณการซื้อขาย คุณสมบัติเชิงเครือข่าย เช่น การกระจุกตัวของคู่เทรด (partner concentration) จึงเป็นสัญญาณสำคัญของพฤติกรรมผิดปกติ

\sect{2.3 Algorithm 1 (Sirolly et al., 2025)}{s23}
Algorithm 1 เป็นวิธีให้คะแนนเชิงเครือข่ายที่คำนวณคะแนนความน่าสงสัย x โดยมีค่าในช่วง [0, 1] ให้แต่ละกระเป๋า แบ่งเป็นสองส่วน
\begin{dotlist}
\item \textbf{ส่วนที่ 1 (คะแนนเริ่มต้น x0):} สัดส่วนปริมาณการซื้อขายของกระเป๋าที่อยู่ในตลาดซึ่งกระเป๋านั้นเคย ``ปิดสถานะ'' (closure) อย่างน้อยหนึ่งครั้ง
\item \textbf{ส่วนที่ 2 (การแพร่คะแนน):} ปรับคะแนนด้วยคู่เทรดตามสูตรวนซ้ำ ½(x0 + Bx) จนลู่เข้า โดย B คือเมทริกซ์สัดส่วนปริมาณการซื้อขายที่กระทำกับแต่ละคู่เทรด
\end{dotlist}
กระเป๋าจะถูกติดธงเมื่อคะแนน x ≥ \gtheta{} และรายการซื้อขายจะถูกนับเป็น wash เมื่อกระเป๋าทั้งสองฝั่งผ่านเกณฑ์ งานวิจัยกำหนดพารามิเตอร์หลัก ได้แก่ เกณฑ์การปิดสถานะ c = 0.005 และใช้ \gtheta\ = 0.90 เป็นค่าหลัก

\sect{2.4 เครื่องมือและไลบรารีที่เกี่ยวข้อง}{s24}
\begin{center}
\begin{tabular}{|L{2.1in}|L{4.0in}|}
\hline
\rowcolor{hdrgray}\textbf{เครื่องมือ} & \textbf{บทบาทในงาน}\\\hline
DuckDB & ประมวลผล SQL บนไฟล์ Parquet หลักสิบล้านแถว (การเตรียมข้อมูล การตรวจปิดสถานะ การรวมยอด)\\\hline
scipy.sparse & เก็บเมทริกซ์คู่เทรด B แบบ sparse และคูณเมทริกซ์-เวกเตอร์ในลูปของอัลกอริทึม\\\hline
NumPy & คำนวณเวกเตอร์คะแนน การลู่เข้า และตัวชี้วัด (recall, precision, AP)\\\hline
pandas & เชื่อมข้อมูลระหว่าง DuckDB กับ NumPy และสร้างตารางสรุป\\\hline
matplotlib & สร้างกราฟประกอบผลการศึกษา\\\hline
uv & จัดการ dependency และสภาพแวดล้อมให้ทำซ้ำได้ (reproducible)\\\hline
\end{tabular}
\end{center}

% =====================================================================
% บทที่ 3
% =====================================================================
\chap{3}{วิธีดำเนินการศึกษา}{ch3}
ระบบออกแบบเป็นสายท่อประมวลผล (pipeline) หลายขั้นที่ส่งต่อกันด้วยไฟล์ ทำให้แต่ละขั้นรันซ้ำและตรวจสอบแยกกันได้ ดังแสดงลำดับต่อไปนี้ (1) ข้อมูลดิบ (2) การเตรียมข้อมูล (3) การให้คะแนน (4) การประเมินผล และ (5) การวิเคราะห์ครบช่วงและการสำรวจข้อมูล

\sect{3.1 แหล่งข้อมูลและการเก็บรวบรวม}{s31}
ใช้ข้อมูลทุติยภูมิ (secondary source) จากชุดข้อมูลสาธารณะ moose-code/polymarket-onchain-v1 บน Hugging Face ขนาดประมาณ 5.3 กิกะไบต์ ประกอบด้วยเหตุการณ์การจับคู่คำสั่งซื้อขาย (order\_\allowbreak filled) ข้อมูลตลาด (condition, market\_\allowbreak data) ข้อมูลถูกจัดเก็บเป็นไฟล์ Parquet รายเดือน

\sect{3.2 การเตรียมข้อมูล (Pre-processing)}{s32}
ขั้นตอนนี้แปลง ``เหตุการณ์การจับคู่ดิบ (raw fill)'' ให้เป็น ``รายการซื้อขายที่จับคู่แล้ว (matched trade)'' หนึ่งแถวต่อหนึ่งรายการ ซึ่งเป็นส่วนที่ท้าทายที่สุด เพราะข้อมูล on-chain มีกับดักหลายจุดที่ต้องจัดการด้วย SQL บน DuckDB ได้แก่
\begin{numlist}
\item ตัดแถวสรุปฝั่งผู้รับคำสั่ง (taker summary) ที่ซ้ำกับแถวฝั่งผู้วางคำสั่ง (maker) โดยดูจากที่อยู่สัญญา exchange
\item แปลงหน่วยจำนวนเงินและจำนวนหุ้นจากหน่วยทศนิยม 6 ตำแหน่ง (หารด้วย 1,000,000)
\item แปลงที่อยู่กระเป๋าให้เป็นตัวพิมพ์เล็กทั้งหมด เพื่อไม่ให้ที่อยู่เดียวกันถูกนับแยก
\item จับคู่โทเคนกับตลาดและผลลัพธ์ (Yes/No) ผ่านตาราง condition.\allowbreak positionIds
\item คำนวณการเปลี่ยนแปลงสถานะของฝั่ง maker โดยอาศัยข้อเท็จจริงว่าการเปลี่ยนสถานะของ taker เท่ากับค่าลบของ maker เสมอ จึงไม่ต้องสร้างฝั่ง taker ใหม่
\item กรองรายการขยะ ได้แก่ การเทรดกับตัวเอง (self-trade) จำนวนหุ้นไม่เป็นบวก และโทเคนที่จับคู่ตลาดไม่ได้
\end{numlist}
ผลลัพธ์คือไฟล์ matched trade พร้อมคอลัมน์ธง in\_\allowbreak target ที่แยก ``สัปดาห์เป้าหมาย'' ออกจาก ``ช่วงประวัติย้อนหลัง (lead-in)'' ซึ่งใช้เพียงเพื่อเรียนรู้สถานะและคะแนนเท่านั้น ตัวชี้วัดหลักคือ share volume (จำนวนหุ้น) ไม่ใช่จำนวนเงิน

\sect{3.3 การพัฒนา Algorithm 1}{s33}
\subsect{3.3.1 การตรวจจับการปิดสถานะ (closure) ด้วย SQL}
สร้างบัญชีเดินรายการ (ledger) ที่ทุกการเทรดปรากฏสองแถว (ฝั่ง maker และ taker) แล้วใช้ window function ไล่สะสมสถานะสุทธิหลังแต่ละเทรด และหาค่าสูงสุดของขนาดสถานะก่อนหน้า จากนั้นติดธงว่ากระเป๋า ``เคยปิดสถานะ'' เมื่อเกิดการกลับทิศของสถานะ หรือการหดตัวครั้งสุดท้ายจนเหลือไม่เกิน 0.5\% ของค่าสูงสุดที่เคยถือ

\subsect{3.3.2 การให้คะแนนและการลู่เข้า}
คำนวณคะแนนเริ่มต้น x0 จากสัดส่วนปริมาณการซื้อขายในตลาดที่พบการปิดสถานะ จากนั้นสร้างเมทริกซ์คู่เทรด B แบบ sparse แล้ววนสูตร ½(x0 + Bx) จนกระทั่งการเปลี่ยนแปลงสัมพัทธ์ (relative L2) น้อยกว่า 1 × 10\textsuperscript{-5} ตรรกะการให้คะแนนถูกรวมไว้ในฟังก์ชันเดียว (score\_\allowbreak wallets) เพื่อให้ทุกสคริปต์ใช้ตรรกะเดียวกัน

\sect{3.4 การทดสอบสมมติฐานและการประเมินผล}{s34}
\begin{dotlist}
\item \textbf{Premise test:} รันอัลกอริทึมกับสอง baseline window แล้วเปรียบเทียบสัดส่วนปริมาณที่ถูกติดธงระหว่างช่วง wash สูงกับต่ำ
\item \textbf{Synthetic injection:} ฉีดวงแหวน wash ที่รู้สมาชิกแน่นอนเข้าไปในช่วงที่เงียบ แล้ววัด recall, precision@k และ average precision (AP) โดยวงแหวนเป็นองค์ประกอบแยกขาดจากกระเป๋าจริง จึงไม่กระทบคะแนนของกระเป๋าจริง
\item \textbf{การตรวจสอบความถูกต้อง (validation):} ตรวจความสมบูรณ์ของชุดข้อมูล และเทียบสถานะที่สร้างใหม่กับบัญชีเดินรายการสำรองที่สร้างจาก raw fill โดยตรง
\item \textbf{Sensitivity analysis:} ทดลองเปลี่ยนค่าเกณฑ์ \gtheta{} และความยาวประวัติย้อนหลัง (30/90/180 วัน) เพื่อดูความคงทนของผล
\end{dotlist}

\sect{3.5 การวิเคราะห์ครบ 18 สัปดาห์และการสำรวจข้อมูล (EDA)}{s35}
รันการให้คะแนนครบ 18 สัปดาห์โดยแต่ละสัปดาห์ใช้ประวัติย้อนหลัง 30 วันของตนเอง การรวมปริมาณข้ามสัปดาห์นับเฉพาะส่วนเป้าหมายเพื่อไม่ให้ประวัติที่ทับกันถูกนับซ้ำ จากนั้นทำการสำรวจข้อมูลเปรียบเทียบกระเป๋าที่ถูกติดธงกับที่ไม่ถูกติดธง พร้อมสร้างกราฟสรุปด้านความไวต่อเกณฑ์ ปริมาณรายวัน การกระจุกตัวของคู่เทรด และสัดส่วนหมวดตลาด

% =====================================================================
% บทที่ 4
% =====================================================================
\chap{4}{ผลการศึกษา}{ch4}
\sect{4.1 ผลการทดสอบสมมติฐาน (Baseline)}{s41}
เมื่อรัน Algorithm 1 ที่ \gtheta\ = 0.90 กับสองสัปดาห์เป้าหมาย พบความแตกต่างของสัดส่วนปริมาณที่ถูกติดธงอย่างชัดเจน ดังตาราง

\begin{center}
\begin{tabular}{|L{3.1in}|r|r|}
\hline
\rowcolor{hdrgray}\textbf{ตัวชี้วัดของสัปดาห์เป้าหมาย} & \textbf{ธ.ค. 2567} & \textbf{มิ.ย. 2568}\\\hline
จำนวนรายการ (trades) & 1,635,178 & 1,147,205\\\hline
กระเป๋าที่ใช้งาน & 104,186 & 87,130\\\hline
จำนวนตลาด & 2,307 & 3,918\\\hline
Share volume & 692.9 M & 389.8 M\\\hline
สัดส่วนที่ถูกติดธง (\gtheta\ = 0.90) & 61.82\% & 0.35\%\\\hline
\end{tabular}
\end{center}

\sect{4.2 ผลการวิเคราะห์ครบ 18 สัปดาห์}{s42}
ผลรวมทั้งสองช่วง (ประมาณ 29.38 ล้านรายการ) ยืนยันความแตกต่างในภาพใหญ่ สัดส่วนปริมาณที่ถูกติดธงเท่ากับ 49.00\% ในช่วง พ.ย.–ธ.ค. 2567 เทียบกับ 2.50\% ในช่วง พ.ค.–มิ.ย. 2568

\begin{center}
\begin{tabular}{|L{3.0in}|r|r|}
\hline
\rowcolor{hdrgray}\textbf{ตัวชี้วัด} & \textbf{พ.ย.–ธ.ค. 2567} & \textbf{พ.ค.–มิ.ย. 2568}\\\hline
Target trades & 19,111,285 & 10,269,237\\\hline
Share volume & 6,712.06 M & 2,667.56 M\\\hline
Flagged share volume & 3,289.07 M & 66.82 M\\\hline
สัดส่วนที่ถูกติดธง (\gtheta\ = 0.90) & 49.00\% & 2.50\%\\\hline
สัดส่วน USDC ที่ถูกติดธง & 39.63\% & 2.14\%\\\hline
ช่วงสัดส่วนรายสัปดาห์ & 13.65–63.76\% & 0.15–5.25\%\\\hline
\end{tabular}
\end{center}

เมื่อจัดตามเดือนตามวันเทรดจริง สัดส่วนที่ถูกติดธงเป็น พ.ย. 2567 = 37.69\%, ธ.ค. 2567 = 61.09\%, พ.ค. 2568 = 4.37\% และ มิ.ย. 2568 = 0.64\% แสดงให้เห็นว่าไม่ควรใช้สัปดาห์ baseline เดียวแทนผลของทั้งสองเดือน

\sect{4.3 ความไวต่อค่าเกณฑ์ (\gtheta-sensitivity)}{s43}
ความแตกต่างระหว่างสองช่วงคงอยู่ที่ทุกค่าเกณฑ์ที่ทดลอง แต่ระดับสัดส่วนเปลี่ยนตามเกณฑ์อย่างมีนัย จึงต้องรายงานควบคู่กับค่าหลักเสมอ

\begin{center}
\begin{tabular}{|c|r|r|}
\hline
\rowcolor{hdrgray}\textbf{\gtheta} & \textbf{สัดส่วน พ.ย.–ธ.ค.} & \textbf{สัดส่วน พ.ค.–มิ.ย.}\\\hline
0.70 & 66.50\% & 18.83\%\\\hline
0.80 & 59.94\% & 7.51\%\\\hline
0.90 & 49.00\% & 2.50\%\\\hline
0.95 & 44.82\% & 1.13\%\\\hline
0.99 & 39.46\% & 0.87\%\\\hline
\end{tabular}
\end{center}

\sect{4.4 การสำรวจข้อมูลกระเป๋าที่ถูกติดธง}{s44}
กระเป๋าที่ถูกติดธงมีปริมาณกระจุกกับคู่เทรดมากกว่ากระเป๋าที่ไม่ถูกติดธงอย่างเห็นได้ชัด เช่น สัดส่วนปริมาณที่กระทำกับคู่เทรดสองรายแรกของกระเป๋าติดธงอยู่ที่ประมาณ 55–57\% เทียบกับ 28–30\% ของกระเป๋าที่ไม่ถูกติดธง อย่างไรก็ตาม เนื่องจากคะแนนใช้ความสัมพันธ์คู่เทรดอยู่แล้ว ผลนี้จึงเป็นการอธิบายการทำงานของตัวตรวจจับ ไม่ใช่การยืนยันอิสระ

\sect{4.5 การประเมินด้วยข้อมูลที่รู้คำตอบ (Synthetic Evaluation)}{s45}
ฉีดวงแหวน wash รวม 144 กระเป๋าใน 6 ระดับความเข้มเข้าในช่วงที่เงียบ ที่ \gtheta\ = 0.90 ตัวตรวจจับจับได้ 48 จาก 144 กระเป๋า (recall = 33.33\%) โดยจับได้เฉพาะวงแหวนที่ความเข้มสูง (≥ 0.95) ส่วนค่า precision@144 = 0 และ average precision (AP) = 0.0046 ซึ่งเป็นค่าต่ำ เนื่องจากกระเป๋าจริงที่ไม่มีเฉลยถูกนับเป็น negative ทั้งหมด คณะผู้จัดทำรายงานค่าต่ำนี้ตามจริงเพื่อความโปร่งใส

\sect{4.6 กรณีที่ต้องตีความอย่างระมัดระวัง}{s46}
\begin{dotlist}
\item \textbf{ตลาด Trump inauguration:} ถูกติดธงสูงถึง 89\% ที่ \gtheta\ = 0.90 แต่ลดเหลือราว 21\% ที่ \gtheta\ = 0.99 จึงเป็นผู้ต้องสงสัยว่าอาจเป็น false positive ต้องตรวจเพิ่มเติม
\item \textbf{ตลาด Lewis Hamilton:} ในสัปดาห์สุดท้ายของเดือน มิ.ย. ปริมาณที่ถูกติดธงเกือบทั้งหมดมาจากตลาดนี้ตลาดเดียว สะท้อนการกระจุกเฉพาะตลาด ไม่ใช่การกลับมาของ wash ในวงกว้าง
\end{dotlist}

% =====================================================================
% บทที่ 5
% =====================================================================
\chap{5}{สรุปผลและข้อเสนอแนะ}{ch5}
\sect{5.1 สรุปผลการศึกษา}{s51}
งานนี้พัฒนาวิธีให้คะแนนเชิงเครือข่ายเพื่อตรวจจับรูปแบบ wash trading บน Poly\-market โดยนำ Algorithm 1 มาพัฒนาใหม่และวิเคราะห์ข้อมูลครบ 18 สัปดาห์ ผลพบความแตกต่างของสัดส่วนปริมาณที่ถูกติดธงอย่างชัดเจน (49.00\% เทียบกับ 2.50\% ที่ \gtheta\ = 0.90) และความแตกต่างนี้คงอยู่ที่ทุกค่าเกณฑ์ที่ทดลอง ซึ่งสนับสนุนว่าตัวตรวจจับสามารถแยกช่วงที่มี wash สูงออกจากช่วงที่ต่ำได้

\sect{5.2 ข้อจำกัด}{s52}
\begin{numlist}
\item ไม่มีข้อมูลเฉลยที่แท้จริง จึงยังไม่ทราบ precision และอัตรา false positive ที่แท้จริง และคะแนนยังไม่ถูกปรับเทียบเป็นความน่าจะเป็น
\item ไม่ทราบสถานะก่อนช่วงประวัติย้อนหลัง สถานะจึงเป็นการประมาณจากศูนย์ ไม่ใช่ยอดคงเหลือจริงในอดีต
\item คะแนนไวต่อค่าเกณฑ์ \gtheta{} และความยาวประวัติ จึงต้องรายงานความไวควบคู่เสมอ
\item การจัดหมวดตลาดใช้คำสำคัญ (keyword heuristic) มีชื่อตกหล่น และการกระจุกคู่เทรดสัมพันธ์กับนิยามคะแนน จึงไม่ใช่หลักฐานอิสระ
\item ผลเป็นการเปรียบเทียบเชิงสังเกตของช่วงเวลาที่เลือก ไม่ได้พิสูจน์สาเหตุหรือเจตนาสมรู้ร่วมคิด
\end{numlist}

\sect{5.3 ข้อเสนอแนะ}{s53}
\begin{numlist}
\item ปรับเทียบคะแนน (calibration) และหาเกณฑ์ที่เหมาะสมเมื่อมีข้อมูลเฉลยบางส่วน
\item ปรับปรุงการจัดหมวดตลาดให้เป็นระบบ และตรวจสอบตลาดที่น่าสงสัยด้วยมือเพื่อยืนยันผล
\item ขยายการประเมินด้วยรูปแบบ wash ที่ซับซ้อนและปะปนกับคู่เทรดจริงมากขึ้น เพื่อทดสอบความทนทานของตัวตรวจจับ
\end{numlist}

% =====================================================================
% บรรณานุกรม
% =====================================================================
\plainchap{บรรณานุกรม}{refs}
\begin{list}{}{\setlength{\leftmargin}{0.5in}\setlength{\itemindent}{-0.5in}\setlength{\itemsep}{0.4em}}
\item Sirolly, A., et al. (2025). Network-based detection of wash trading on prediction markets. (Algorithm 1, \S5.1 และ \S6).
\item moose-code. (2025). polymarket-onchain-v1 [ชุดข้อมูล]. Hugging Face. CC-BY-4.0. \url{https://huggingface.co/datasets/moose-code/polymarket-onchain-v1}
\item DuckDB Foundation. (2024). DuckDB: An in-process SQL OLAP database management system. \url{https://duckdb.org}
\end{list}

\end{document}
